\section*{Chapter \ref{ch:rc:vanilla-rates-options} --- Vanilla Rates Options}

\begin{solution}{exo:rc:vanilla-rates-options:1}
$10\,000\,000\times(0.034-0.030)\times0.5069 = \text{EUR}~20\,276$, paid at the
end of the period, six months after the fixing.
\end{solution}

\begin{solution}{exo:rc:vanilla-rates-options:2}
$\sigma_B\approx\sigma_N/F = 80/250 = 32\%$; the exact value 32.7\% is slightly
higher because $2\Phi(x/2)-1 < x/\sqrt{2\pi}$ for $x = \sigma_B\sqrt T > 0$.
\end{solution}

\begin{solution}{exo:rc:vanilla-rates-options:3}
The first period's rate is already fixed on the trade date: its caplet has no
optionality left, only a known payoff, so it is left out and the cap's
optionality starts with the second period.
\end{solution}

\begin{solution}{exo:rc:vanilla-rates-options:4}
Above: the two-year cap's value at 66 basis points includes the first-year
caplets at 58; for the average to reach 66 the second-year caplets must be
higher. The strip gives 66.6.
\end{solution}

\begin{solution}{exo:rc:vanilla-rates-options:5}
The physical price multiplies the option by the collateralised annuity 7.740;
the par-yield formula uses $P(0,T)a(S_0) = 7.579$, which discounts the
annuity at the Euribor swap rate 3.098\%, higher than the €STR discount rates:
2.1\% less. The collateralised cash price settles at the swap's value on the
clearing house's \omterm{def:rc:multi-curve-and-collateral-discounting:curves}{discount curve}, which is what the physical swaption delivers,
so both prices coincide.
\end{solution}

\begin{solution}{exo:rc:vanilla-rates-options:6}
In the normal model a move of so many basis points is equally likely at any
level; relative to a low strike that is a large percentage move, relative to a
high strike a small one, so the implied \omterm{def:rc:vanilla-rates-options:lognormal}{lognormal volatility} falls with the
strike. A flat lognormal market looks like a positive normal skew: normal
volatility $\approx\sigma_B\sqrt{FK}$ rises with the strike.
\end{solution}

\begin{solution}{exo:rc:vanilla-rates-options:7}
1.97\%: selling the five-year floor at 1.97\% pays for the cap at 3\%. The
company then pays at least 1.97\% plus its margin even if Euribor falls
further. (Using the 3\% caplet volatilities at every strike ignores the smile.)
\end{solution}

\begin{solution}{exo:rc:vanilla-rates-options:8}
76 basis points is the \omterm{def:rc:vanilla-rates-options:flat}{flat volatility} of the whole five-year cap, an average
over all its caplets; the caplets fixing after four years have their own
stripped volatility, 78.0. At 76 they are worth EUR~886\,956 on EUR~200 million,
at 78.0 EUR~916\,329: underpriced by EUR~29\,373.
\end{solution}

\begin{solution}{pb:rc:vanilla-rates-options:1}
\textbf{1.} Nine: ten six-month periods, less the first, already fixed.
\textbf{2.} EUR~2\,157\,482, 1.079\% of the notional.
\textbf{3.} Divided by the annuity of the caplet payments (4.286): 25.2 basis
points a year.
\textbf{4.} The same, EUR~2\,157\,482: the stripped volatilities were calibrated so
that the five-year cap reprices at its \omterm{def:rc:vanilla-rates-options:flat}{flat volatility}.
\textbf{5.} The forwards are expectations; with 66 to 79 basis points of normal
volatility a year, fixings above 3\% in the later years are quite likely, and
the cap pays in those states.
\textbf{6.} EUR~43\,093 per basis point.
\textbf{7.} About $10\times43\,093 = \text{EUR}~430\,929$.
\textbf{8.} Grow: one point of \omterm{def:rc:vanilla-rates-options:lognormal}{lognormal volatility} is worth about $F$ points of
normal volatility ($\sigma_N\approx F\sigma_B$), so the lognormal vega is about
$F$ times the normal vega, which does not change with the level; higher rates,
larger lognormal vega. Normal vega is the more stable number to report.
\textbf{9.} The bank that sold it; it hedges by buying caps or caplets from the
market, or swaptions for the longer expiries, and runs the residual.
\textbf{10.} EUR~1\,827\,377: at a flat 25\% the caplets on forwards around
2.4\% have normal volatilities of about 60 basis points, below the market's;
the lognormal quote would misprice the cap by 15\%.
\textbf{11.} $200\text{ million}\times0.005\times\sum\delta_k$ with $\sum\delta_k =
4.569$: EUR~4\,569\,444.
\textbf{12.} EUR~4\,285\,872 ($0.005\times200\text{ million}\times$ annuity 4.286).
\textbf{13.} About 3.25\%: the premium is 25.2 basis points a year of
the notional, so the fixings must average 25 basis points above the strike.
\textbf{14.} Insurance: protection against the states in which rates rose,
which did not occur; the premium is the price of that protection.
\textbf{15.} Each caplet pays on its own fixing, so a high fixing in one period is
not offset by a low one in another; an option on the average would be cheaper.
\textbf{16.} The swap fixes the rate at the forward (about 2.4\% to 3\% over the
periods) at no upfront cost and removes the benefit of falling rates; the cap
costs 25 basis points a year and keeps it.
\textbf{17.} 1.97\%: the company pays for its cap by giving up the benefit of
fixings below 1.97\%.
\textbf{18.} Bid--offer on the volatility it must buy to hedge, the credit charge
for a corporate that does not post collateral (chapters~18 and 19), and the
capital the trade consumes.
\textbf{19.} Named result: \emph{the treasurer's cap} costs EUR~2\,157\,482 (1.079\%
of notional, 25.2 basis points a year) with a vega of EUR~43\,093 per basis
point of normal volatility.
\textbf{20.} Protection against high rates while keeping the benefit of low ones.
\end{solution}

\begin{solution}{iq:rc:vanilla-rates-options:1}
An option paying $N\delta\max(F-K,0)$ at the end of a period on the rate fixed at
its start. Its forward $F_k(t) = (P(t,T_{k-1})/P(t,T_k)-1)/\delta$ is a traded
asset divided by $P(t,T_k)$, so it is a martingale under the $T_k$-forward
measure, which makes the caplet price $N\delta P(0,T_k)$ times a call on a
martingale.
\iqlookfor{the numeraire argument.}
\end{solution}

\begin{solution}{iq:rc:vanilla-rates-options:2}
At the money, Bachelier gives $\sigma_N\sqrt{T}\varphi(0) = \sigma_N\sqrt{T/2\pi}$
and Black gives $F(2\Phi(\sigma_B\sqrt T/2)-1)\approx F\sigma_B\sqrt{T}\varphi(0)$.
Equating: $\sigma_N\approx F\sigma_B$; more generally $\sigma_N\approx\sigma_B\sqrt{FK}$
for nearby strikes.
\iqlookfor{the expansion and the $\sqrt{FK}$ rule.}
\end{solution}

\begin{solution}{iq:rc:vanilla-rates-options:3}
Price each quoted cap at its \omterm{def:rc:vanilla-rates-options:flat}{flat volatility}; with caps in increasing maturity,
solve for one volatility of the new caplets so that the longer cap reprices,
earlier caplets held (or fit a smooth parametric curve to all caps at once).
Problems: negative or oscillating volatilities when quotes are inconsistent or
maturities close, sensitivity of the strip to one bad quote, a different strip
for every strike (the smile), and the first-caplet convention.
\iqlookfor{the sequential solve and its failure modes.}
\end{solution}

\begin{solution}{iq:rc:vanilla-rates-options:4}
Normal volatility is more stable when rates move, since rate changes in basis
points depend little on the level; it works with negative rates; normal vegas
add across instruments of different forwards; and the market quotes it, so hedge
ratios match the quotes.
\iqlookfor{stability of the smile and additivity.}
\end{solution}

\begin{solution}{iq:rc:vanilla-rates-options:5}
The par-yield payoff $a(S_T)(S_T-K)^+$ is not the annuity numeraire times the
payoff, so annuity-measure pricing does not apply; the formula
$P(0,T)a(S_0)\,\mathrm{Black}$ ignores the convexity of $a(S)$ and mixes
discount rates, and it is not arbitrage-free against physical swaptions. The
market (euro swaptions, November 2018) moved the default to the collateralised
cash price, which settles at the swap's value on the clearing house's curve.
\iqlookfor{payoff versus numeraire, and the market fix.}
\end{solution}

\begin{solution}{iq:rc:vanilla-rates-options:6}
Black's formula takes $\ln(F/K)$: a negative forward or strike gives a NaN or
an exception. Price in the normal model or a shifted lognormal model (chapter
5), convert the volatility quotes accordingly, and test the pricer on negative
forwards and strikes.
\iqlookfor{the log and the model switch.}
\end{solution}
